\documentclass[11pt]{article}
\usepackage{graphicx} % Required for inserting images
\usepackage{newtxtext,newtxmath}
\usepackage{geometry}
\geometry{a4paper,top=25mm,}

\title{WASP COURSE SOFTWARE ENGINEERING ASSIGNMENT 2026}
\author{Zelong Wang\\Linköping University\\zelong.wang@liu.se}
\date{Aug. 2026}

\begin{document}

\maketitle

\section{Introduction}
My research focuses on remote ID for commercial drones and leveraging remote ID for collision avoidance. With the rapidly growing number of commercial drones in everyday life, it is necessary and required (by the Federal Aviation Administration (FAA) and European Union Aviation Safety Agency (EASA)) to assign a digital id to every civil drone, so the public (can be everyone with a smartphone) know that every operating drone is legal and safe (like every car has a plate). Therefore, aviation authorities proposed the Remote ID standard, which requires every drone in flight to continuously broadcast its unique digital ID and real-time spatial information such as velocity, GNSS position, and the operator's location; therefore, any nearby observer can verify the drone using a smartphone.

However, the most adopted standard proposed by the American Society for Testing and Materials (ASTM) specifies no security protection: the remote ID broadcast is plain text, leaving it vulnerable to attacks such as spoofing and tampering. My work is to enable security measures such as cryptography and physical layer security in the remote ID to make it more secure and robust, so every public receiver knows that the transmitted data from the drone is correct and truthful~\cite{survey_rid}. The implemented secure remote ID protocol is called the Drone Remote ID Protocol (DRIP)~\cite{card2022drone,formal_verification}, which was initially proposed by the Internet Engineering Task Force (IETF). 

Furthermore, I will leverage the secure remote ID information for further drone management applications. Nearby drones could use the message to know the precise location of other drones and then avoid colliding with each other~\cite{trustbased_ca}. I will use development kits and small prototype drones, such as Nrf kits and Crazyflie, as the testbed to implement the remote ID technique and collision avoidance algorithm. The goal of the research is to improve public safety in the air and increase the efficiency of civic drone traffic management. To sum up, the research consists of three main phases: building a secure remote ID architecture with enhanced security measures for civil drones on the original non-secure remote ID; enabling collision avoidance based on remote ID information; and implementing the architecture and algorithm on a real-world testbed.


\section{Lecture principles}
I want to discuss AI/ML Technical Debt first, which is particularly important for me~\cite{technical_debt}. My research does not focus on the AI-based system itself, but generative AI can support my work and make it more efficient. I need to run many simulations and real-life tests for my research; therefore, I have to work with many hardware and SDK platforms. From the lecture, I understand that technical debt exists throughout the SE life cycle, which is important for my experiments. 
On one hand, generative AI could create new technical debt for SE because LLMs are based on probability~\cite{technical_debt}. However, it can also reduce human-introduced technical debt in some ways. \cite{technical_debt} found that it is hard to say GenAI creates more debt than human-written code; instead, GenAI creates debt in different ways from humans. What matters to me is how to take advantage of GenAI while minimizing technical debt. For instance, when using the Nrf SDK, I need to learn the device tree concept and lots of technical documents; AI could make this learning curve much easier and optimize the standard device tree structure. Whereas Gen-AI may create high-risk security vulnerabilities, it is not the most important concern for an early-phase demo, and it can be addressed once the demo/research has been successfully tested. 

Secondly, I would like to discuss "AI Changes Developer Behavior" and psychology from a human aspect, which can be considered as a follow-up idea of the first concept. The work in~\cite{AIcode_insecure} shows that every AI coding assistant introduces at least one issue; more importantly, it found that, in the comparison groups, programmers using AI assistants report higher confidence in their code. Since their code can pass the test run, they tend to believe the quality of AI assistant code and ignore potential security risks. This psychological condition could cause technical debt to snowball and introduce long-term maintenance costs in the SE life cycle. Vibe coding, or AI assistant coding, could accelerate programming significantly and change developer behavior, but it also shifts more effort to code review. Another paper compared human-written code and AI-generated code~\cite{2025humanwrittenvsaigeneratedcode}; a key finding is that it is hard to say whether human or AI introduces more technical debt. Instead, human and AI introduce issues in different ways. Human-written code shows higher logic complexity and lower readability, but it usually contains fewer security risks. On the contrary, AI-generated code tends to be more compact and better structured to achieve the functionality, but it may contain potentially high-risk vulnerabilities. These papers offer insightful guidance on how to interact with LLMs more safely. I have been working with generative AI and vibe coding to understand its strengths and weaknesses. This will help me use it for more specific tasks in controlled environments, maximizing its advantages and minimizing security risks.

\section{Guest-lecture principles}
I think requirement engineering is an important topic from the guest lecture. This topic is not directly related to my research, as my work does not focus on AI or software engineering. Nevertheless, with the rapid development of LLMs and generative AI, I have been working with it more often. The lecture introduced a clear SE life cycle for students without an SE background, and more specifically, the focus on requirements engineering and how it could integrate into AI-assisted workflows. Knowing the problem is more important than knowing the solutions, because if we don't know the exact problem, how should we come up with a proper solution? The logic of RE can be integrated into any ; requirements can vary across different areas, but the methodology remains the same. The concept of RE can be easily adopted to most fields. 

Secondly, I would like to discuss the importance of human aspects in SE. As AI becomes more integrated into SE workflows, developers may tend to work more independently. This can make limited human interaction more critical; otherwise, insufficient communication may make the whole process hard to synchronize and consolidate. My research involves wide cooperation with many other teams from different countries and cultures, such as France, Brazil, and the IETF working group from North America. Each group is responsible for a different module of the big project. Hence, the human aspect is particularly important for our remote collaboration, such as how we interact and communicate with each other and how we synchronize our work. My research can be considered a semi-SE life-cycle, as our workflow includes processes from requirements to real-world prototype and implementation. This topic can help my research group improve overall teamwork and performance.


\section{Data scientists, software engineers, and AI engineers}
Data scientists usually have an educational background and focus on model building and data pre-processing. They tend to evaluate their work in terms of accuracy and robustness, rather than qualities such as interference latency and training costs. Whereas software engineers tend to deliver software products within a given budget that meet the user's needs. They often work with uncertainty and budget constraints, apply engineering judgment to handle trade-offs between various qualities, and ultimately deliver a software product for a customer. I generally agree with the distinction, although it is a simplified classification. This is because educational and professional backgrounds can strongly influence how people think and work, and the career requirements for scientists and engineers differ, which can shape their behaviors. A system-wide approach is critical in my research area; my research includes a comprehensive framework for drone security that requires researchers and engineers from different countries and backgrounds to work togetheso understanding eachaother's others' work is important for understanding the whole pipeltheir own and theiitta scientists and engineers both need to learn about others' skills to build an interdisciplinary team; team members from different backgrounds need to understand each other. AI engineering aligns closely with system engineering for AI systems, as it focuses on the infrastructure for machine learning. However, the term tends to be model-centric rather than system-wide. Whereas data engineering usually focuses on data management and infrastructure. Lastly, foundation models, AI coding assistants, etc., could change these boundaries because these techniques bridge the gap between scientists and engineers by eliminating information barriers. Since agentic AI can handle more concrete tasks, humans are expected to have a global perspective to cover all aspects of a large system. 

\section{Paper analysis}
This section provides the analysis of two full papers from CAIN, I decide to choose \textit{DDPT: Diffusion-Driven Prompt Tuning for Large Language Model Code Generation}~\cite{DDPT} from year 2025, and \textit{Cognition Envelopes for Bounded Decision Making
in Autonomous UAS Operations}~\cite{2026cognitionenvelopesboundeddecision} from year 2026. The reason is that the first paper discusses optimizing the prompt tuning for AI code generation that is very timely and practical for current work flow for both academia and industry. The second paper address the use of LLMs and VLMs for decision making of drones, the focus of cyber-physical system and UAS operations align well with my own research. 

\subsection{Paper 1: DDPT: Diffusion-Driven Prompt Tuning for Large Language Model Code Generation}
The paper introduces a novel automated prompt engineering framework, Diffusion-Driven Prompt Tuning (DDPT), for using large language models (LLMs) in code generation. DDPT utilizes a diffusion model to generate optimal continuous prompt embeddings from Gaussian noise rather than using fixed prompt parameters. In this process, the LLM remains frozen. The diffusion model is updated based on the code-generation loss to learn the optimal distribution of prompt embeddings. From an engineering perspective, the proposed method bypasses the massive computational overhead normally required to fine-tune large foundation models by optimizing embeddings directly. The DDPT model automates the traditional manual process of prompt engineering. The evaluation indicates that this approach significantly improves the quality of AI-generated code. It improves both the syntactic structure and the semantic accuracy of the generated code and provides a scalable, efficient method to enhance AI performance in the software development life cycle.

This paper does not directly relate to my research topic, but it could contribute to my workflow from an application perspective. At the moment, I work with a lot of video coding and AI-generated code to make demos and prototypes; therefore, I work a lot with prompt engineering for generated code. Based on my experience, a major issue with current AI-generated code is translating human language into AI-understandable language accurately; prompts I think are correct can still produce unexpected outputs. Often, people have to revise the prompt and the output back and forth. Hence, this paper provides a practical method to improve the vibe coding workflow and potentially solve a critical problem. As long as the LLM can understand human prompts more accurately, it can reply with more accurate output and increase overall workflow efficiency, producing generated code that is functionally and syntactically more accurate. 

For a fictional AI-intensive software project, this project features a coding system that directly translates natural language into software modules. Its main functionality is translating user intent into functional code. By integrating the concept from this paper, this system could generate more accurate code without additional computational overhead. When a user submits a requirement, DDPT can work in the background to automatically optimize prompt embeddings. This automated prompt optimization could ensure the LLMs receive highly tuned instructions, giving code with improved syntactic and semantic accuracy across diverse programming tasks. As vibe coding has seen growing adoption and is being integrated into a growing number of software projects, the idea of this paper could fit with a large number of projects quite easily. As long as my research involves AI-generated code to some extent, this paper could fit my work and improve code accuracy. 

My project cannot be tweaked toward the idea of the DDPT paper; adopting this paper is more about application and cognition. Nevertheless, one challenge remains unsolved: LLM hallucinations. DDPT could make LLMs understand better, but the AI model can still produce hallucinated outputs. This problem is not solved in the paper, so keep it in mind when working with vibe coding and review the code carefully before implementation. 

\subsection{Paper 2: Cognition Envelopes for Bounded Decision Making
in Autonomous UAS Operations}
This paper introduces the concept of cognition envelopes, a runtime assurance layer that uses a reasoning boundary to constrain decisions made by LLMs and vision-language models (VLMs). Unlike traditional safety envelopes that rely on physical boundaries such as speed or geofencing, or meta-cognition methods that inherit the same hallucinations because AI critiques itself, the proposed cognition envelopes can serve as an independent monitor. This method checks the AI decision against external physical evidence, mission rules, objectives, and uncertainty. If a decision violates constraints, a gating rule will revise the decision or escalate it to a human operator. In AI engineering, this concept is vital because foundational AI models can introduce new vulnerabilities, such as hallucinations and over-generalizations, that may lead to defective decision-making. In cyber-physical systems such as autonomous systems or drone management, safety and security are critical; flawed decisions may compromise objectives or endanger human lives. The proposed cognition adds an extra guardrail layer to detect improper plans before execution and revise or escalate the decision to a human operator. Therefore, the concept improves the trustworthiness and accountability of the UAS system, enhancing overall safety and security.

The proposed paper relates to my research pretty straightforwardly. Although AI is not my main topic, my research focuses heavily on the safety and security of UAS and cyber-physical systems. My recent research consists of secure communication and collision avoidance. While the current collision avoidance algorithm makes decisions instantly, it may lead to failed operations, creating a gap that could be addressed by adding a secure layer to revise the avoidance decision based on the actual physical situation. There is great potential to integrate this decision guardrail to enhance the overall security and effectiveness of collision avoidance.

Considering a larger AI project for UAS management, the project is to integrate an AI system with drones. All drones are equipped with a collision avoidance system that makes decisions based on incoming communication from other drones; therefore, the cognition envelopes are responsible for handling complex real-time physical environments and revising decisions from the collision avoidance system. For instance, when a drone receives a broadcast from two other drones, the collision avoidance trigger asks for an instant maneuver, but such a plan might not be proper for the current task or situation regarding the physical surroundings; in this case, the cognition envelope layer could re-evaluate the plan based on external evidence and decide to revise the plan or escalate it to the human operator. 

This paper can be integrated with my research pretty directly. As my research focuses on the security of drones and cyber-physical systems, the proposed paper aligns very well with my topic. By employing secure communication, a drone could receive broadcasts from other drones and perform collision-avoidance maneuvers or collaborate. Currently, popular decision-making approaches include AI-based and non-AI-based algorithms. For both types of methods, a filter gate to revise the decision is still missing. Implementing the cognition envelopes from this paper in current decision-making systems could improve decision effectiveness and reduce flight risks, as discussed above.

\section{Evidence, originality, and GenAI-use transparency}
All the main resources I have used are in the reference list. Sometimes I may have a glimpse of another resource, such as Google, news, a business report, or a wiki, but I use it mainly for a quick impression or background on topics, not for formal citations.

I sometimes use GenAI for brainstorming and search help. If I have an idea for a topic, I can ask AI for different perspectives or ask it to challenge my idea, such as with a red team test. When GenAI gives me arguments or opinions, I will ask for links or references to support them. By doing this, I can find the original source for a manual check. If I find a concrete source like a peer-reviewed paper, I quickly review it and think about the argument in more detail. If the source is a fake link, I ignore such a reply. Lastly, I also use AI (Grammarly) for language checking, which could correct typos and grammar mistakes while keeping the original human-written text as much as possible. 

The sources I trust mostly are peer-reviewed papers and IETF drafts, which I use most in my research. Web pages or blogs can also be helpful; I use them mainly for informational purposes, such as getting a brief background on an unfamiliar area or catching the most recent news. The reason is simple: a peer-reviewed paper requires experts in the same field to approve the work, which sets a quality bottom line. In contrast, anyone can publish a product page or personal blog at any time. So the average quality bar for a peer-reviewed paper is higher than for other free sources.


\bibliographystyle{IEEEtran}
\bibliography{biblio}

\end{document}
